\section{Performance}

\label{sec:performance}

All benchmarks on AMD EPYC 7763, single core, Go~1.22 with AVX-512 SIMD NTT.

\subsection{TFHE Boolean Operations}

\begin{table}[ht]
\centering
\caption{TFHE gate evaluation benchmarks ($N = 1024$, 128-bit security).}
\label{tab:tfhe-benchmarks}
\begin{tabular}{lr}
\toprule
\textbf{Operation} & \textbf{Latency} \\
\midrule
AND / OR / XOR / NAND / NOR / XNOR & 12\,ms \\
NOT & 0.02\,ms \\
MUX (conditional select) & 24\,ms \\
8-bit addition & 310\,ms \\
8-bit comparison & 280\,ms \\
16-bit addition & 680\,ms \\
16-bit comparison & 620\,ms \\
\bottomrule
\end{tabular}
\end{table}

\subsection{BFV Batched Arithmetic}

\begin{table}[ht]
\centering
\caption{BFV batched arithmetic benchmarks ($N = 4096$, 128-bit security).}
\label{tab:bfv-benchmarks}
\begin{tabular}{lrr}
\toprule
\textbf{Operation} & \textbf{16-bit} & \textbf{32-bit} \\
\midrule
Encrypted addition & 45\,ms & 52\,ms \\
Encrypted subtraction & 45\,ms & 52\,ms \\
Encrypted multiplication & 180\,ms & 340\,ms \\
Plaintext-ciphertext multiply & 22\,ms & 28\,ms \\
Relinearization & 85\,ms & 120\,ms \\
Bootstrapping & 1.2\,s & 2.4\,s \\
\bottomrule
\end{tabular}
\end{table}

\subsection{Comparison to Cleartext EVM Execution}

\begin{table}[ht]
\centering
\caption{FHE overhead relative to cleartext EVM operations.}
\label{tab:overhead}
\begin{tabular}{lrrr}
\toprule
\textbf{Operation} & \textbf{Cleartext} & \textbf{FHE (BFV)} & \textbf{Overhead} \\
\midrule
\texttt{uint256 add} & 3 gas (5\,ns) & 65K gas (45\,ms) & $9 \times 10^6$ \\
\texttt{uint256 mul} & 5 gas (8\,ns) & 320K gas (180\,ms) & $2.3 \times 10^7$ \\
\texttt{>=} comparison & 3 gas (5\,ns) & 480K gas (480\,ms) & $9.6 \times 10^7$ \\
\bottomrule
\end{tabular}
\end{table}

The overhead is large in absolute terms ($10^6$--$10^8\times$). This is inherent to
FHE and consistent across all implementations. The relevant comparison is not FHE
vs.\ cleartext, but FHE vs.\ the \emph{alternative for achieving the same privacy
guarantee}:

\begin{itemize}[leftmargin=2em]
  \item Off-chain computation with on-chain verification (ZK proofs): requires a
        trusted prover and does not support arbitrary computation.
  \item Trusted execution environments (TEEs): requires trusting hardware vendors
        and is vulnerable to side-channel attacks.
  \item Multi-party computation (MPC): requires interaction between parties for
        every operation and does not scale to public smart contracts.
\end{itemize}

FHE is the only approach that enables \emph{non-interactive} computation on encrypted
data by \emph{untrusted} parties. The overhead is the price of this unique property.

%% ========================================================================
%%  8. SECURITY ANALYSIS
%% ========================================================================
